\section{Counting excitation subsets}

For $0\le\tau\le1/2$, write
$h(\tau)=-\tau\log\tau-(1-\tau)\log(1-\tau)$, with $0\log0=0$.
All logarithms are natural.

\begin{theorem}[Binary-entropy bound for binomial coefficients]
    \label{thm:excitation_binomial_entropy}
    \lean{Real.log_choose_le_mul_binEntropy, Real.choose_le_exp_mul_binEntropy}
    \leanok
    Let $k,r$ be nonnegative integers and $0\le\tau\le1/2$.
    If $r\le\tau k$, then
    \begin{align}
        \log\binom{k}{r}&\le kh(\tau),\\
        \binom{k}{r}&\le e^{kh(\tau)}.
    \end{align}
    These are the binomial estimates used before the inverse-compression
    inequality in~\cite[Section~7]{OpenAI2026AreaLaw}.
\end{theorem}

\begin{proof}\leanok
    If $\tau=0$, then $r=0$ and both claims are equalities.
    Otherwise $r\le k$ and each nonnegative term in the binomial expansion
    of $(\tau+(1-\tau))^k=1$ is at most $1$. Hence
    \begin{align}
        \log\binom{k}{r}
        &\le-r\log\tau-(k-r)\log(1-\tau)\\
        &=kh(\tau)+(\tau k-r)\log\frac{\tau}{1-\tau}
        \le kh(\tau).
    \end{align}
    The last inequality follows from $\tau\le1/2$ and $r\le\tau k$.
    Exponentiating proves the second claim.
\end{proof}

\begin{theorem}[Binary-entropy bound for excitation subsets]
    \label{thm:excitation_subset_entropy}
    \lean{Real.sum_choose_le_mul_exp_binEntropy,
        Finset.card_filter_powerset_le_mul_exp_binEntropy}
    \leanok
    \uses{thm:excitation_binomial_entropy}
    Let $S$ be a finite set with $k$ elements and $0\le\tau\le1/2$.
    Then
    \begin{align}
        \sum_{\substack{0\le r\le k\\r\le\tau k}}\binom{k}{r}
        &\le(k+1)e^{kh(\tau)},\\
        \bigl|\{B\subseteq S:|B|\le\tau k\}\bigr|
        &\le(k+1)e^{kh(\tau)}.
    \end{align}
    This is the subset-count estimate preceding the inverse-compression
    inequality in~\cite[Section~7]{OpenAI2026AreaLaw}.
\end{theorem}

\begin{proof}\leanok
    There are at most $k+1$ admissible cardinalities, each contributing at
    most $e^{kh(\tau)}$ by Theorem~\ref{thm:excitation_binomial_entropy}.
    Every admissible subset belongs to the family of subsets of cardinality
    $r$ for one admissible $r$, and that family has $\binom{k}{r}$ members.
\end{proof}
